\begin{abstract}
We address a single scope-locked question: what are the formal limits of
KPI-based enterprise analysis when enterprises are modeled as continua
\(K_{EA}\) under the Ontology of Continua (OC) and its executable theory
specification \texttt{ETS.K\_EA.v1.1}.
Within this framework, KPIs are treated as admissible readout maps defined on
the OC/ETS observable domain
\(\Omega_{\mathrm{obs}}\subseteq\Omega(K_{EA})\), where observability itself is a
structural property of the enterprise continuum, comprising an axis-level
observability component, a factor \(S_{\mathrm{obs}}\), and an
ObservabilityInvariant.

KPI-based observation is formalized as a quotienting of
\(\Omega_{\mathrm{obs}}\) by the fibers of the KPI readout map.
On this basis, we establish existence and genericity results showing that
KPI observation is intrinsically non-identifying: structurally distinct
enterprise states may remain indistinguishable under any finite admissible KPI
family, regardless of sampling resolution, aggregation schemes, or downstream
deterministic statistical post-processing.

These results are strictly diagnostic.
They delineate fundamental information-theoretic limits of KPI-driven analysis
and characterize what cannot be recovered even in principle in
continuum-modeled enterprise systems.
\end{abstract}

\noindent\textbf{Keywords:}
enterprise continuum; KPI observability; identifiability limits;
non-identifiability; impossibility results
